% Source-derived chapter 06: The case of two nef isotropic classes
% Original source: riemann-roch-polynomials-hyperkahler-fourfolds
\chapter{The case of two nef isotropic classes}
\label{riemann-roch-polynomials-hyperkahler-fourfolds-chapter-06}
\source{riemann-roch-polynomials-hyperkahler-fourfolds}

\begin{remark}[Source section]
\label{riemann-roch-polynomials-hyperkahler-fourfolds-chapter-06-source}
\source{riemann-roch-polynomials-hyperkahler-fourfolds}
\source{riemann-roch-polynomials-hyperkahler-fourfolds}
This chapter reproduces the corresponding section of the retrieved
LaTeX source, including its definitions, statements, examples, and
proofs. It has not yet been translated into Lean.
\notready
\end{remark}

% The source section title was: The case of two nef isotropic classes
\label{sec:TwoNefIsotropic}
\source{riemann-roch-polynomials-hyperkahler-fourfolds}

Let $(X,\lll,\mm)$ be a very general triple satisfying~\eqref{hypo}.\
By Proposition~\ref{prop:SurfacesOnHK4}, we have  $\NS(X)=\Z\lll\oplus\Z\mm$.\
We assume in this section that we are in case~\ref{caseC1}: the class $\mm$ is isotropic and both~$\lll$ and $\mm$ are nef.\ Our aim is to prove Proposition~\ref{prop:CaseC1} (in Section~\ref{proof}).

As observed in Section~\ref{subsec:HKNefCone}, all cones of divisors are equal and the class $p \lll+q \mm$ is ample on $X$ for all integers $p,q>0$.\
By Kodaira vanishing, we get
\begin{equation}\label{eq:KodairaVanishingTwoNef}
\source{riemann-roch-polynomials-hyperkahler-fourfolds}
h^0(X,L^{\otp p}\otimes M^{\otp q})=\chi(X,  L^{\otp p}\otimes M^{\otp q})=P_{RR,X}(2pq) =\binom{pq+3}{2}
\end{equation}
and in particular
\[
h^0(X,L\otimes M)=6,\quad h^0(X, L^{\otp 2}\otimes M)=10,\quad h^0(X, L^{\otp 3}\otimes  M^{\otp 2})=36.
\]
The next two sections and the paper  \cite{voisinfib} are devoted to the  study of  the induced rational map
\[
\phi_{L\otimes M}\colon X\dra \P^5.
\]
This map was studied by O'Grady in~\cite[Section~4]{og1} when~$X$ is   of~K3$^{[2]}$ numerical type, and for a very general deformation of the pair $(X,L\otimes M)$.\ 


\subsection{Case where $L$ and $M$ both have nonzero sections}\label{nonz}
\source{riemann-roch-polynomials-hyperkahler-fourfolds}

The following lemma will allow  us to apply Lemma~\ref{lem:H0large}.\ We keep the same hypotheses: $(X,\lll,\mm)$ is a very   general triple satisfying~\eqref{hypo} and we are in case~\ref{caseC1}.

\begin{lemm}\label{lecasholhoM}
\source{riemann-roch-polynomials-hyperkahler-fourfolds}
Assume that $H^0(X,L)\ne 0$ and $H^0(X,M)\ne 0$.\
Then
\begin{equation*}
h^0(X,L)+h^0(X,M)\geq 3,
\end{equation*}
hence either $h^0(X,L)\geq 2$ or $h^0(X,M)\geq2$.
\end{lemm}

\begin{proof}
Let $\sigma$ be a nonzero section of $L$ and let $\tau$ be a nonzero section of $M$.\
The divisors $D_\sigma$ of $\sigma$ and $D_\tau$ of $\tau$   have no common component, since the class of any effective divisor on $X$ is an integral combination of $\lll$ and $\mm$ with nonnegative coefficients.\
It follows that
\[
\Sigma_{\sigma\tau}\coloneqq D_\sigma\cap D_\tau
\]
is a surface whose ideal sheaf $\cI_{\Sigma_{\sigma\tau}}$ admits the Koszul resolution
\begin{equation*}\label{eqkosresSigma}
\source{riemann-roch-polynomials-hyperkahler-fourfolds}
0\to  (L\otimes M)^{-1}\to    L^{-1}\oplus  M^{-1}\to   \cI_{\Sigma_{\sigma\tau}} \to  0.
\end{equation*}
If we tensor it by $L\otimes M$, the associated  exact sequence in cohomology gives
\begin{equation}\label{eqineq1}
\source{riemann-roch-polynomials-hyperkahler-fourfolds}
h^0(X,\cI_{\Sigma_{\sigma\tau}}(L\otimes M))=h^0(X,L)+h^0(X,M)-1.
\end{equation}
If we tensor the resolution by $L^{\otp 2}\otimes M^{\otp 2}$, we get, using~\eqref{eq:KodairaVanishingTwoNef},
\begin{equation}\label{eqineq2}
\source{riemann-roch-polynomials-hyperkahler-fourfolds}
\begin{split}
h^0(X,\cI_{\Sigma_{\sigma\tau}}(L^{\otp 2}\otimes M^{\otp 2}))&=h^0(X,L^{\otp 2}\otimes M)+h^0(X,L\otimes M^{\otp 2})-h^0(X,L\otimes M )=14,\\ h^1(X,\cI_{\Sigma_{\sigma\tau}}(L^{\otp 2}\otimes M^{\otp 2}))&=0.
\end{split}
\end{equation}
Using again~\eqref{eq:KodairaVanishingTwoNef}, we get $h^0(X,L^{\otp 2}\otimes M^{\otp 2})=21$, and, from~\eqref{eqineq2}, we deduce
\begin{equation}\label{eqineq3}
\source{riemann-roch-polynomials-hyperkahler-fourfolds}
h^0(\Sigma_{\sigma\tau},(L^{\otp 2}\otimes M^{\otp 2})\vert_{\Sigma_{\sigma\tau}})=7.
\end{equation}

Assume by contradiction $h^0(X,L)+h^0(X,M)=2$.\
Then, by~\eqref{eqineq1}, we get
\begin{equation}\label{eqineq4}
\source{riemann-roch-polynomials-hyperkahler-fourfolds}
\rk \bigl(H^0(X, L\otimes M)\lra  H^0(\Sigma_{\sigma\tau}, ( L\otimes M)\vert_{\Sigma_{\sigma\tau}})\bigr)= 5.
\end{equation}
In particular, the surface $\Sigma_{\sigma\tau}$ is not contained in the base locus of $ |L\otimes M|$.\

We prove now the following properties.

\begin{claim}
\source{riemann-roch-polynomials-hyperkahler-fourfolds}
\notready\label{claimavecc2}
\source{riemann-roch-polynomials-hyperkahler-fourfolds}
In the notation used above,
\begin{enumerate}[{\rm (a)}]
\item\label{enum:Sigma1} the surface $\Sigma_{\sigma\tau}$ is irreducible and reduced;
\source{riemann-roch-polynomials-hyperkahler-fourfolds}
\item\label{enum:Sigma2} the image of the rational map
\source{riemann-roch-polynomials-hyperkahler-fourfolds}
\[
\phi_{L\otimes M}\vert_{\Sigma_{\sigma\tau}}\colon  \Sigma_{\sigma\tau}\dra \P^5
\]
is a surface.
\end{enumerate}
\end{claim}


\begin{proof}
To prove~\ref{enum:Sigma1}, we note that the surface $\Sigma_{\sigma\tau}$ has class $\lll\mm$ and its associated matrix (as in~\eqref{eqmatriceO}) is
\[
M_{[\Sigma_{\sigma\tau}]}=\begin{pmatrix}
0 & 2 \\
2 & 0
\end{pmatrix}.
\]
If $\Sigma_{\sigma\tau}$ is not irreducible or not reduced, there   exist surfaces $\Sigma_1$, $\Sigma_2$ in $X$, such that
\begin{equation*}
M_{[\Sigma_1]}=M_{[\Sigma_2]}=\begin{pmatrix}
0 & 1 \\
1 & 0
\end{pmatrix}.
\end{equation*}
We show that no such   decomposition  exists under our assumptions.\

Since $(X,\lll,\mm)$ is very general, we may apply Proposition~\ref{prop:SurfacesOnHK4} and write
\[
[\Sigma_i]=t_{i} \lll^2+u_{i}\lll\mm+v_{i} \mm^2+w_{i} q_X^\vee
\]
with $t_{i}, u_i,v_i\in \Q_{\ge0} $ (because $\lll$ and $\mm$ are nef) and $w_i\in \Q $.\
Using \eqref{eqdematl2m2}, we find
\begin{equation*}\label{eqpoureta3eta4}
\source{riemann-roch-polynomials-hyperkahler-fourfolds}
t_{i}=v_i=0,\quad 2u_{i}+25w_{i}=1,
\end{equation*}
so that
\[
[\Sigma_i]=\frac{1}{2}\Bigl(1-25w_{i}\Bigr) \lll\mm+w_{i} q_X^\vee.
\]
We obtain $w_2=- w_1$, so we can assume $w=:w_1\geq 0$.\
We have
\begin{equation*}\label{eqSigmai}
\source{riemann-roch-polynomials-hyperkahler-fourfolds}
[\Sigma_i]=\frac{1}{2}\lll\mm-(-1)^i w\Bigl(q_X^\vee-\frac{25}2\lll\mm\Bigr),\quad \textnormal{for }i\in\{1,2\}.
\end{equation*}
Using~\eqref{qxc}, we get
\begin{equation*}\label{eqnumbres}
\source{riemann-roch-polynomials-hyperkahler-fourfolds}
\Sigma_1^2= \Sigma_2^2=\frac{1}{2}+525w^2,\qquad \Sigma_1\Sigma_2=\frac{1}{2}-525w^2.
\end{equation*}
Since these numbers are integers, $w$ is positive and $2\cdot w^2=1\cdot 3\cdot 7\cdot (5w)^2$ is an integer, which implies that the denominator of $w$ is at most $5$, so that $w\geq \frac{1}{5}$.\
Then,
\begin{equation}\label{eqeffectivec2}
\source{riemann-roch-polynomials-hyperkahler-fourfolds}
w q_X^\vee=[\Sigma_1]+\frac{1}{2}(25w -1)\lll\mm
\end{equation}
with $25w -1>0$.\ But this is impossible: by Proposition~\ref{lem:NefCone4folds} and~\cite[Proposition 3.2]{huyKahler}, the closure of the K\"ahler cone of $X$ is the subset $\overline{\cK}_X\subset H^{1,1}(X,\R)$ consisting of those real $(1,1)$-classes $\omega$ satisfying $q_X(\omega)\geq 0$, $q_X(\omega,\omega_0)>0$ for a K\"ahler class $\omega_0$, and $q_X(\omega,\lll)\ge0$ and $ q_X(\omega,\mm)\geq0$.
For all $\omega\in \overline{\cK}_X$ and all effective classes $\eta\in{\Hdg}^4(X,\Q)$, we have $\int_X\eta\omega^2\geq 0$.
Furthermore, for all $\omega\in \overline{\cK}_X$ such that $q_X(\omega)=0$, we have $\int_X q_X^\vee \omega^2=0$; so by~\eqref{eqeffectivec2} we have
\[
0=\int_{\Sigma_1}\omega^2+\frac{1}{2}(25w -1)\int_X\lll\mm\omega^2,
\]
where both terms on the right are nonnegative.\
Hence $\int_X\lll\mm\omega^2=0$; this implies that the class~$\lll\mm$ is proportional to $q_X^\vee$, which is absurd.

Claim~\ref{enum:Sigma2}, namely the fact that the image $\phi_{L\otimes M}(\Sigma_{\sigma\tau})$ is a surface, is proved as follows.\
If $\phi_{L\otimes M}(\Sigma_{\sigma\tau})$ is a linearly nondegenerate curve in $\P^4$, it has degree~$d\ge4$.\ Let $\fff$ be the class of a fiber of $\Sigma_{\sigma\tau}\dra\phi_{L\otimes M}(\Sigma_{\sigma\tau})$.\
 We have
\[
(\lll+\mm)[\Sigma_{\sigma\tau}]=(\lll+\mm)\lll\mm = \lll^2\mm + \lll\mm^2
\]
and this class can be written as $d\fff+\eee$, where $\eee$ is an effective curve class.\  This contradicts the fact that the class of any effective curve in $X$ has the form $t \lll^2\mm+u \lll\mm^2$  with   $t, u\in\frac12\Z_{\ge0}$.\end{proof}
The claim gives us a contradiction, since we know by~\eqref{eqineq4} that $\phi_{L\otimes M}(\Sigma_{\sigma\tau})$ is a linearly nondegenerate surface in $\P^4$, while, by~\eqref{eqineq3}, it is contained in at least~$\binom{4+2}{2}-7=8$
independent quadrics; this contradicts Castelnuovo's lemma, which says that it is contained in at most $\binom{2+1}{2}$ independent quadrics.
\end{proof}


\subsection{Case where either $L$ or $M$   has no nonzero sections}\label{zero}
\source{riemann-roch-polynomials-hyperkahler-fourfolds}
We keep the same hypotheses:~$(X,\lll,\mm)$ is a very   general triple satisfying~\eqref{hypo} and we are in case~\ref{caseC1}.\ This section is devoted to the proof of the following result.

\begin{prop}
\source{riemann-roch-polynomials-hyperkahler-fourfolds}
\notready\label{prop:CaseB'2}
\source{riemann-roch-polynomials-hyperkahler-fourfolds}
If  either $H^0(X,L)=0$ or $H^0(X,M)=0$, the image of the rational map $\phi_{L\otimes M}\colon X\dra \P^5$ is rationally connected.
\end{prop}

Our aim is to prove that
 \[
Y\coloneqq\Im(\phi_{L\otimes M})\subset \P^5  
\]
is rationally connected, that is, that a general pair of points of $Y$ are joined by a rational curve.\ 
If \mbox{$\dim (Y)<4$,} this holds  by \cite[Theorem~1.4]{linthese}, so it suffices to prove Proposition~\ref{prop:CaseB'2} in  the case $\dim (Y)=4$, that is, when $Y\subset \P^5$ is a hypersurface.\

 We will assume  $H^0(X,L)=0$; the case  $H^0(X,M)=0$ is of course analogous (just permute~$L$ and $M$).\ 

\begin{prop}
\source{riemann-roch-polynomials-hyperkahler-fourfolds}
\notready\label{prop6.3precise} If   $H^0(X,L)=0$ and $\dim (Y)=4$, a plane section $C_{0,W_3}$ of $Y$ defined by a general vector subspace $W_3\subset H^0(X,L\otimes M)$ of dimension $3$   is either  of geometric genus~$0$ or  a smooth cubic curve.
\source{riemann-roch-polynomials-hyperkahler-fourfolds}
\end{prop}

Proposition \ref{prop6.3precise} implies Proposition~\ref{prop:CaseB'2} as follows.\  Proposition~\ref{prop6.3precise} then  shows that either~$Y$ is  rationally connected since  its general plane sections are rational curves, or it is  a cubic hypersurface, hence it is uniruled.\ In the second case, we can   apply \cite[Theorem~1.4]{linthese} again to the maximal rationally connected quotient of $Y$, which has dimension $<4$.\ It is thus rationally connected and is therefore a point by~\cite{GHS}, proving again that $Y$ is rationally connected.

We   start the proof of Proposition~\ref{prop6.3precise} by introducing some notation.\
Set
\begin{align*}
&W_6\coloneqq H^0(X,L\otimes M),\\
&W_{10}\coloneqq H^0(X,L^{\otp 2}\otimes M),\\
&W_{36}\coloneqq H^0(X,L^{\otp 3}\otimes  M^{\otp 2}),
\end{align*}
with multiplication map
\[
\mu\colon W_6\otimes W_{10}\lra W_{36}.
\]

\begin{lemm}\label{lem:NoRk2}
\source{riemann-roch-polynomials-hyperkahler-fourfolds}
If $H^0(X,L)=0$,  there are no rank-2 elements in $\Ker(\mu)$.
\end{lemm}

\begin{proof}
A rank-$2$ element in~$\Ker(\mu)$ is given by a nontrivial relation
$\alpha\sigma=\beta \tau $ in $H^0(X,L^{\otp 3}\otimes M^{\otp 2})$ with $\alpha,\,\beta\in H^0(X,L\otimes M)$ linearly independent.\ Since $H^0(X,L)=0$, any divisor in $|L\otimes M|$ is irreducible and reduced, hence the divisors of $\alpha$ and $\beta$  have no common component.\ It follows that $\tau\in H^0(X,L^2\otimes M)$ vanishes on the divisor of $\alpha$ hence can be written as the product of~$\alpha$ by a section of $L$.\ This implies $\tau=0$, which is absurd.
\end{proof}

\begin{lemm}\label{leW6W10}
\source{riemann-roch-polynomials-hyperkahler-fourfolds}
Assume there are no rank-$2$ elements in $\Ker(\mu)$.\
For a general $3$-dimensional vector subspace $W_3\subset W_6$, the restriction $\mu\vert_{W_3}\colon W_3\otimes W_{10}\to W_{36}$ of the map $\mu$ has rank $\geq 28$.
\end{lemm}

\begin{proof}
Let $\cS_3\to \Gr(3,W_6)$ be the rank-$3$ tautological subbundle.\
The natural sheaf inclusion $\cS_3\to  W_6\otimes \cO_{\Gr(3,W_6)}$ induces a morphism $f\colon \P(\cS_3\otimes W_{10})\to\P(W_6\otimes W_{10})$ which makes $\P(\cS_3\otimes W_{10})$ a smooth birational model of the set of elements of rank at most $ 3$ in $\P(W_6\otimes W_{10})$.\
We set
\[
R_3\coloneqq f^{-1}(\P(\Ker(\mu))).
\]
Since, by assumption, there are no rank-$2$ elements in $\Ker(\mu)$, the scheme $R_3$ is isomorphic to the set of elements of rank $\leq 3$ in $\P(\Ker(\mu))$.

The fiber of the natural morphism $\pi\colon R_3\to\Gr(3,W_6)$ over a point $[W_3]$ is the space $\P(\Ker(\mu\vert_{W_3}))$.\
Arguing by contradiction, if the conclusion of the lemma does not hold, the fiber of $\pi$ has dimension $\geq 2$, hence $\dim (R_3)\geq 11$.

Let $H$ be  the restriction to $R_3$ of the line bundle $\cO_{\P(\cS_3\otimes W_{10})}(1)$.\
We have a tautological inclusion $H^{-1}\to \pi^* \cS_3\otimes W_{10}$
\begin{equation}\label{eqmorphdual}
\source{riemann-roch-polynomials-hyperkahler-fourfolds}
W_{10}^\vee\otimes \cO_{R_3}\lra  H\otimes \pi^*\cS_3.
\end{equation}
Since there are no rank-$2$ elements in $\Ker (\mu)$,  the morphism~\eqref{eqmorphdual}  is surjective, with kernel $\cK $ a locally free sheaf  of rank $7$.\
Thus we have $c_i(\cK)=0$ for $i>7$, that is, $s_i(H\otimes \pi^*\cS_3)=0$, for all $i>7$, where the $s_i$ denote the Segre classes.


We now deduce a contradiction.\
For any line bundle $H$ on a variety and any vector bundle~$\cE$ of rank $3$, we have the relation
\begin{eqnarray}\label{eqrel}
\source{riemann-roch-polynomials-hyperkahler-fourfolds}
s_i(\cE\otimes H)=\sum_{j=0}^i (-1)^j \binom{i+2}{j} H^js_{i-j}(\cE).
\end{eqnarray}
In our case, one has $s_j(\cS_3)=0$ for $j\geq 4$, because of the exact sequence
\[
0\to   \cS_3\to  W_6\otimes \cO_{\Gr(3,W_6)}\to  \cQ_3\to  0
\]
on the Grassmannian, where $\cQ_3$ is locally free of rank $3$.\
The relations \eqref{eqrel}  thus give four linear relations involving  $s_0(\pi^* \cS_3),\ldots, s_3( \pi^* \cS_3)$, namely
\begin{align*}\label{eqrelexpl}
\source{riemann-roch-polynomials-hyperkahler-fourfolds}
  0&=s_8(\pi^* \cS_3\otimes H  )\nonumber\\
  &= \binom{10}{8}H^8s_0(\pi^* \cS_3)-\binom{10}{7}H^7s_1(\pi^* \cS_3)+\binom{10}{6}H^6 s_2(\pi^* \cS_3)-\binom{10}{5}H^5s_3(\pi^* \cS_3),\\
  \nonumber
 0&=s_9(\pi^* \cS_3\otimes H  )\\
 &= -  \binom{11}{9}H^9s_0(\pi^* \cS_3)+\binom{11}{8}H^8s_1(\pi^* \cS_3)-\binom{11}{7}H^7s_2(\pi^* \cS_3)+\binom{11}{6}H^6s_3(\pi^* \cS_3), \\
  \nonumber
 0&=s_{10} (\pi^* \cS_3\otimes H  )\\
 &=\binom{12}{10}H^{10}s_0(\pi^* \cS_3)-\binom{12}{9}H^9s_1(\pi^* \cS_3)+\binom{12}{8}H^8s_2(\pi^* \cS_3)-\binom{12}{7}H^7s_3(\pi^* \cS_3), \\
  \nonumber
 0&=s_{11} (\pi^* \cS_3\otimes H  )\\
 &=-\binom{13}{11}H^{11}s_0(\pi^* \cS_3)+\binom{13}{10}H^{10}s_1(\pi^* \cS_3)-\binom{13}{9}H^9s_2(\pi^* \cS_3)+\binom{13}{8}H^8s_3(\pi^* \cS_3).
\end{align*}

We can assume $\dim(R_3)=11$, replacing it by a proper algebraic subset if necessary.\
Multiplying these equations by   adequate powers of $H$, we get the linear relations
\begin{equation*} 
\begin{split}
& 0=H^3s_8(\pi^* \cS_3\otimes H),\qquad 0=H^2s_9(\pi^* \cS_3\otimes H  ),\\
& 0=H s_{10} (\pi^* \cS_3\otimes H  ),\qquad 0=s_{11} (\pi^* \cS_3\otimes H),
\end{split}
\end{equation*}
which, by expanding as above,  give four linear relations between the four intersection numbers
\begin{equation*} \label{eqnumbers}
\source{riemann-roch-polynomials-hyperkahler-fourfolds}
H^8s_3(\pi^* \cS_3),\ H^9s_2(\pi^* \cS_3),\ H^{10}s_1(\pi^* \cS_3), \ H^{11}
\end{equation*}
on $R_3$.\
Since these four  linear relations are clearly independent, we conclude that these four intersections numbers vanish, which contradicts the fact that $H$ is ample on $R_3$.
\end{proof}

Lemma~\ref{lem:NoRk2} and Lemma~\ref{leW6W10} together imply the following result.

\begin{coro}\label{lealternative}
\source{riemann-roch-polynomials-hyperkahler-fourfolds}
Assume $H^0(X,L)=0$.\ For a general $3$-dimensional vector subspace $W_3\subset W_6$, the image $\mu(W_3\ot  W_{10})$ has dimension $\geq 28$.
\end{coro}

\begin{proof}[Proof of Proposition~\ref{prop6.3precise}.]
We choose a general $3$-dimensional vector subspace $W_3\subset W_6$.\
The locus
defined by the vanishing of the sections in $W_3$ has a mobile part $C_{W_3}\subset X$ which, by Bertini's theorem, is an irreducible curve
which   dominates the  plane curve $C_{0,W_3}$ via $\phi_{L\otimes M}$.\ Consider the restriction maps
\[
r^{p,q}\colon H^0(X,L^{\otp p}\otimes M^{\otp q})\lra  H^0(C_{W_3},(L^{\otp p}\otimes M^{\otp q})\vert_{C_{W_3}})
\]
for $p,q>0$.\
Let us set $W_{p,q}'\coloneqq \Im(r^{p,q})$.\
We will estimate the dimension of $W_{p,q}'$, for $(p,q)\in\{(1,1),(2,1), (3,2)\}$.\
First of all, we note that $W_{1,1}'$ has dimension~3.

We then claim that $W_{3,2}'$ has dimension $\leq 8$.
Indeed, we have the inclusion
\[
\mu(W_3\ot W_{10})\subset \Ker (r^{3,2})
\]
and, by Corollary~\ref{lealternative},   the space on the left has dimension at least $28$, while $h^0(X,L^{\otp 3}\otimes M^{\otp 2})=36$.

Finally, we claim that
\begin{enumerate}[{\rm (a)}]
\item\label{enum:claim1_part1} either the space $W_{2,1}'$ has dimension $\geq 5$,
\source{riemann-roch-polynomials-hyperkahler-fourfolds}
\item\label{enum:claim1_part2} or it has dimension $4$ and the image $\phi_{L^{\otp 2}\otimes M}( C_{W_3})$ is a rational normal cubic curve in~$\P^3$.
\source{riemann-roch-polynomials-hyperkahler-fourfolds}
\end{enumerate}
Indeed, since $\phi_{L\otimes M}$ is a dominant rational map to a hypersurface $Y\subset \P^5$, the curve $C_{W_3}$ can be chosen to pass through a  general triple of points $x,y,z\in X$.\
Assume  $\phi_{L^{\otp 2}\otimes M}( C_{W_3})$ spans at most a $\P^3$ in $\P^9=\P(W_{10}^\vee)$.\
This $\P^3$ then contains the projective plane  spanned by the images of $x$, $y$,  $z$, and thus this plane intersects the curve $\phi_{L^{\otp 2}\otimes M}(C_{W_3})$ in at least a fourth point, unless   $\phi_{L^{\otp 2}\otimes M}(C_{W_3})$ is a rational normal curve of degree $3$.\

In the former case, we conclude that the variety $Y'\coloneqq\phi_{L^{\otp 2}\otimes M}(X)$, which spans $\P^9$, has the property that a general trisecant plane of $Y'$ is $4$-secant.\
This is absurd, since a general $1$-dimensional linear section of $Y'$ spans at least a $\P^6$.\

In the latter case, the curve $\phi_{L^{\otp 2}\otimes M}(C_{W_3})$ has  degree $\leq 3$.\
This curve cannot be a plane curve, since otherwise the projection of $Y'$ through any of its points $y$ would contain a line through any two points $x$, $z$, hence would be a projective space of dimension at most $4$.\
Hence it must be a rational normal curve in $\P^3$, as claimed.

Let us now consider the multiplication map
\[
W_{1,1}'\otimes W_{2,1}'\lra  W_{3,2}'.
\]
In case~\ref{enum:claim1_part1}, the three spaces have respective dimensions $3$, at least $5$, and at most $8$.\
Since the curve $C_{W_3}$ is irreducible, we can apply Lemma \ref{lepourconcprop72} below: it says that   this is possible only if the linear system $W_{1,1}'$ factors through an elliptic plane  curve  or  a rational curve.

In case~\ref{enum:claim1_part2}, we have $W_{2,1}'=\Sym^3(W_2'')$, for some $2$-dimensional linear system $W_2''$ on~$C_{W_3}$.\
We now consider the multiplication maps
\begin{align*}
W_{1,1}'\otimes W_2''&\lra  W_{k_1}'',\\
W_{k_1}''\otimes W_2''&\lra  W_{k_2}'',\\
W_{k_2}''\otimes W_2''&\lra  W_{3,2}',
\end{align*}
where $W_{k_1}''$ and $W_{k_2}''$ are spaces of sections of adequate line bundles on~$C_{W_3}$, of respective dimensions $k_1$ and $k_2$.\
The Hopf lemma (\cite[p.~108]{acgh}) gives
\begin{align*}
\dim (W_{k_1}'')&\geq \dim (W_{1,1}')+1,\\
\dim (W_{k_2}'')&\geq \dim (W_{k_1}'')+1,\\
\dim (W_{3,2}')&\geq \dim (W_{k_2}'')+1.
\end{align*}
As $\dim (W_{1,1}')=3$ and $\dim (W_{3,2}')\leq 8$, the three inequalities above cannot all be strict, hence one of them must be an equality.\
It is well known that this implies  that the plane curve $C_{0,W_3}=\phi_{L\otimes M}(C_{W_3}) $ is rational, thus completing the proof of Proposition \ref{prop6.3precise}.
\end{proof}

We used  above the following lemma, for which we could not find a reference.

\begin{lemm}\label{lepourconcprop72}
\source{riemann-roch-polynomials-hyperkahler-fourfolds}
Let $C$ be a smooth connected projective curve and let $H$ and $H'$ be  line bundles on $C$.\ Let
$W_3\subset H^0(C,H)$ and $W_k\subset H^0(C,H')$ be  base-point-free linear systems on $C$, of respective linear dimensions $3$ and $k$, with $k\geq 4$.\
Assume that the rank of the multiplication map
$$\mu\colon W_3\otimes W_k\lra H^0(C,H\otimes H')$$
is at most $3+k$.\ Then, both linear systems factor through a morphism $C\to C_0$, where $C_0$ is either a  rational curve or a degree-$3$ elliptic curve in $\P(W_3)$.
\end{lemm}

\begin{proof} Denote by $\phi_3\colon C\to \P^2$   the morphism  induced by $W_3$ and by
$\phi_k\colon C\to \P^{k-1}$ the morphism  induced by $W_k$.\   We first claim that the result is true if  $\phi_{k}$ does not factor through~$\phi_3$.\ Indeed, assume there is a length-$2$ subscheme
$z\subset C$ that imposes only one  condition on $W_3$ and two conditions on $W_k$.\ Then
$z$ imposes two conditions on $\Im(\mu)$, hence the multiplication map
$$\mu_z\colon W_3(-z)\otimes W_k\lra   H^0(C, H\otimes H')$$
has rank at most $k+1$, while $W_3(-z)$ has dimension $2$.\ We are thus in the  equality case of the Hopf lemma  and we conclude that $W_k$ is the pullback to $C$ of $H^0(\P^1,\cO_{\P^1}(k-1))$ via a morphism
$\psi\colon C\to  \P^1$.\ This case is easily concluded by studying  the
multiplication maps
$$W_3\otimes H^0(\P^1,\cO_{\P^1}(k-2))\lra H^0(C,H\otimes \psi^*\cO_{\P^1}(k-2)),$$
 with image $W'$, and
$$W'\otimes H^0(\P^1,\cO_{\P^1}(1))\lra H^0(C,H\otimes \psi^*\cO_{\P^1}(k-1)), $$
of rank $\leq 3+k$.\
By the Hopf lemma applied to both maps,   $\dim (W')\in\{k+1,k+2\}$.\ Thus, one the two maps  satisfies  the equality in the Hopf lemma, hence    $\phi_3$ also factors through $\psi$.\ This proves the  claim.

We now prove a similar claim with $\phi_3$ and $\phi_{k}$ permuted.\ Assume  there is a length-$2$ subscheme $z\subset C$ that imposes only one condition on $W_k$ but two conditions on $W_3$.\ This produces a $(k-1)$-dimensional vector subspace $W_{k}(-z)\subset W_k$ of sections vanishing on $z$, with the property that the multiplication map
$$\mu_z\colon W_3\otimes W_{k}(-z)\lra H^0(X,H\otimes  H'(-z))$$
has rank $\leq k+1$, which is the minimum allowed by the Hopf Lemma.\ We conclude that there is a morphism $\psi:C\rightarrow \P^1$ such that $W_3=\psi^*H^0(\P^1,\cO_{\P^1}(2))$.\
We now consider the   multiplication maps
$$\mu'_1\colon  W_2\otimes  W_k\lra H^0(C, H'\otimes \psi^* \cO_{\P^1}(1)),$$
with image $W'$, and
$$ \mu'_2\colon W_2\otimes  W'\lra H^0(C, H'\otimes \psi^* \cO_{\P^1}(2))=H^0(C,H\otimes H').$$
We know that  $\mu'_2$ has rank at most $3+k$, so either $\dim (W')=k+1$ and $\mu'_1$ satisfies  the equality case of  the Hopf lemma, or
$\dim (W')=k+2$ and  $\mu'_2$ satisfies  the equality case of  the Hopf lemma.\ In both cases, we conclude that both linear systems factor through $\psi$.

Using these two claims, we can now assume that both linear systems factor through  the curve $C_0\coloneqq\phi_3(C)\subset \P(W_3)$,  that $C_0$ is birationally isomorphic to its image in $\P(W_k)$, and that  the normalization  $C'_0$ of $C_0$ is not rational, as otherwise the lemma is proved.\ We  denote by~$H_0$ and $H'_0$ the line bundles on  $C'_0$  whose respective pullbacks to $C$ are $H$ and $H'$.\      For    general points $x_1,\ldots,x_{k-3}\in C'_0$, we have a
$3$-dimensional space $W'_3
\subset W_k$ of sections vanishing at $x_1,\ldots,x_{k-3}$, and the multiplication map
$$\mu'\colon  W_3\otimes W'_3\lra H^0(C'_0, H_0\otimes H'_0)$$
has rank $\leq 6$ since its image is contained in $ \Im (\mu)$ and vanishes at $x_1,\dots,x_{k-3}$.
We can thus assume that the kernel
\begin{equation}
\label{eqKsubset}
\source{riemann-roch-polynomials-hyperkahler-fourfolds}
K\coloneqq  \Ker (\mu')\subset  W_3\otimes W'_3
\end{equation}
has dimension $3$ (if it has dimension $4$,   we are in the equality case of the Hopf lemma, so we can ignore this case).\
The inclusion \eqref{eqKsubset} induces a morphism
\begin{equation}
\label{eqmorphKW}
\source{riemann-roch-polynomials-hyperkahler-fourfolds}
K\otimes \cO_{\P(  W_3)}\lra  W'_3\otimes  \cO_{\P(  W_3)}(1)
\end{equation}
of rank-$3$ vector bundles
on $\P(W_3)$.\ This morphism  has rank at most $ 2$ along $C_0$ and   must have rank at least $2$ generically on $\P(W_3)$ since otherwise, the image would be  a rank-$1$ sheaf with at least three independent sections contained in $W'_3\otimes  \cO_{\P(  W_3)}(1)$, hence a copy of  $\cO_{\P(  W_3)}(1)$ and the elements of~$K$ would be rank-$1$ tensors in $W_3\otimes W'_3$.\   Since we assumed that the curve $C'_0$ is not rational, no quadratic equation vanishes on it, hence  the morphism \eqref{eqmorphKW} also has generic rank $2$ along the curve $C_0$.

If the curve  $C_0$ has degree at most $ 3$,  the lemma is proved.\ Otherwise, the morphism~\eqref{eqmorphKW}  has   rank at most $ 2$ everywhere on $\P(  W_3)$ and is generically of rank $2$ along $C_0$, which means that the three polynomials of type $(1,1)$ on $\P( W_3)\times \P(  W'_3)$ given by $K$ vanish on  a surface $\Sigma$ which contains the natural  embedding of $C'_0$ in $\P( W_3)\times \P(  W'_3)$ and is birationally isomorphic to  $\P(  W_3)$ by the first projection.\ The surface $\Sigma$  is  also birationally isomorphic  to  $\P( W'_3)$ by the second projection since its image in $\P( W'_3)$ contains the image of $C'_0$ in $\P( W'_3)$ which, being birationally isomorphic to $C'_0$ for a generic choice of $x_1\ldots,x_{k-3}$, is also not rational.\

We claim that the surface $\Sigma\subset \P(  W_3)\times \P( W'_3)$ is  the graph of an isomorphism $\P(  W_3)\isom \P( W'_3)$.\ Indeed, as $\Sigma $ is  contained in three hypersurfaces of type $(1,1)$, it is an irreducible component of the complete intersection of two such hypersurfaces, but it is not the complete intersection of two such hypersurfaces (since there is a third equation of type $(1,1)$ vanishing on it).\  It follows that the class of $\Sigma$ is of the form $(h_1+h_2)^2-e=h_1^2+2h_1h_2+h_2^2-e$, where  $h_i\coloneqq c_1({\rm pr}_i^*\cO_{\P^2}(1))$ and the class $e$  is effective and nonzero on $\P(  W_3)\times \P( W'_3)$.\ Since the   projections ${\rm pr}_1$ and ${\rm pr}_2$ are dominant and $(\Sigma\cdot h_1\cdot h_2)^2\ge (\Sigma\cdot h_1^2)(\Sigma\cdot h_2^2)$ (Hodge Index Theorem), the only possibility is   $[\Sigma]=h_1^2+h_1h_2+h_2^2$, that is, $[\Sigma]$ is the class of the graph of an isomorphism $\P(  W_3)\cong \P( W'_3)$.\ This implies that $\Sigma$ itself is the graph of an isomorphism since we then have $\Sigma^*\cO_{\P( W'_3)}(1)= \cO_{\P( W_3)}(1)$.\ This proves the claim.\
The claim   implies that the
line bundles  $H_0$ and $H'_0(-x_1-\dots-x_{k-3})$ on $C'_0$ coincide.\ As $x_1\ldots,x_{k-3}$ are general points of   $C'_0$ and $k\geq 4$, this implies that $C'_0$ is rational, which is a contradiction.
\end{proof}


\subsection{Proof of Proposition~\ref{prop:CaseC1}}\label{proof}
\source{riemann-roch-polynomials-hyperkahler-fourfolds}
If $H^0(X,L) $ and $H^0(X,M) $ are both nonzero, by Lemma~\ref{lecasholhoM} and Lemma~\ref{lem:H0large}, after possibly permuting $L$ and $M$, the line bundle $L$ is globally generated and thus gives a Lagrangian fibration $f\colon X\to \P^2$ with $f^*\cO_{\P^2}(1)\isom L$.\ So we are in case \ref{prop64a} of Proposition~\ref{prop:CaseC1}.

If either $H^0(X,L) $ or $H^0(X,M) $ is zero, the image of the rational map~$\phi_{L\otimes M}$ is rationally connected by Proposition~\ref{prop:CaseB'2} and we are in case \ref{prop64b} of Proposition~\ref{prop:CaseC1}.\


%%%%%%%%%%%%%%%%%%%%%
